\subsection{Арифметические операции}
Операции: \verb|+ - * / % ^| (плюс, минус, умножение, деление, остаток, возведение в степень). Сравнения: \verb|< <= > >= == ~=|. В Lua 5.1 нет операторов целочисленного деления \verb|//| и возведения \verb|**|; используйте \verb|math.floor(a/b)| и \verb|a ^ b| соответственно. Все числа — \hyperlink{term:number}{вещественные}, поэтому деление отдаёт дробные значения; для округлений применяются \verb|math.floor|, \verb|math.ceil| и \verb|math.abs|.

Приоритет: сначала выполняются \verb|^|, затем \verb|* / %|, далее \verb|+ -|. Скобки меняют порядок вычислений. Остаток \verb|%| определён как «остаток от деления», полезен для задач чётности и циклических индексов.
\begin{minted}{lua}
-- приоритет: умножение выполняется раньше сложения
local a = 2 + 3 * 4
-- остаток от деления (modulo)
local r = 5 % 2
-- возведение 2 в степень 3
local p = 2 ^ 3
-- сравнение: равно?
local eq = (a == 14)
-- сравнение: больше?
local gt = (a > 10)
-- безопасное «целочисленное» деление с округлением вниз
local q = math.floor(17 / 4)  -- 4
-- евклидова длина вектора
local len = math.sqrt(3*3 + 4*4)  -- 5
\end{minted}
\paragraph{Пояснения.}
- \verb|math.sqrt| возвращает корень квадратный; \verb|math.pi| — число $\pi$.
- \verb|math.max| и \verb|math.min| ограничивают значения: \verb|v = math.max(0, math.min(1, v))|.
- Деление на ноль следует избегать: проверяйте знаменатель перед делением.
\subsection*{Задачи для самостоятельного решения}
\begin{enumerate}
  \item Вычислите выражение \verb|3 + 4 * 2 - 5|.
  \item Найдите остаток \verb|17 % 4|.
  \item Запишите \verb|3 ^ 4|.
  \item Вычислите среднее \verb|(x + y) / 2|.
  \item Запишите \verb|math.abs(x2 - x1)|.
  \item Вычислите \verb|math.floor(n / 3)|.
  \item Запишите условие «\verb|n| чётное» с \verb|%|.
  \item Вычислите \verb|math.sqrt(x*x + y*y)| (евклидова длина).
  \item Ограничьте \verb|v| в диапазоне \verb|[0,1]|: \verb|v = math.max(0, math.min(1, v))|.
  \item Переведите градусы в радианы: \verb|rad = deg * math.pi / 180|.
  \item Вычислите площадь прямоугольника \verb|S = w * h|.
  \item Запишите среднее трёх чисел \verb|(a+b+c)/3|.
  \item Реализуйте линейную интерполяцию \verb|lerp = a + t*(b-a)|.
  \item Посчитайте сумму массива \verb|a|: примените цикл \verb|for|.
  \item Найдите максимум массива \verb|a|: пройдитесь и сравнивайте.
\end{enumerate}
